% Chapter 4

\chapter{System Architecture}

\label{Chapter4}

\lhead{Chapter 4. \emph{System Architecture}}

This chapter describes the architecture of the PCIe-AXI DMA streaming pipeline implemented for this thesis: a seven-stage, purely behavioural SystemVerilog model that reproduces the latency behaviour of Equations~\ref{eq:pcie_latency}--\ref{eq:total_latency} while remaining fully synthesisable and cycle-simulatable. \autoref{sec:ch4_overview} gives the pipeline overview and its AXI4-Stream signalling convention; \autoref{sec:ch4_pcie_model}--\ref{sec:ch4_latency_monitor} describe each of the seven stages in turn; \autoref{sec:ch4_top} describes the top-level integration and the parameters exposed for the experiments of \autoref{Chapter6} and \autoref{Chapter7}.

%----------------------------------------------------------------------------------------

\section{Pipeline Overview}
\label{sec:ch4_overview}

The datapath is organised as a linear chain of seven modules, each connected to its neighbours by an AXI4-Stream interface, as shown in \autoref{fig:pipeline_overview}. A packet generator (implemented in the testbench, \autoref{Chapter6}) injects packets into \texttt{pcie\_model}, which are then carried, one stage at a time, through the DMA model, an AXI memory-mapped-to-stream adapter, a parameterised FIFO, an optional batching unit, and finally a stream sink; a hardware latency monitor observes injection and completion events and computes per-packet latency directly from the simulated hardware.

\begin{figure}[htbp]
\centering
\begin{tikzpicture}[
  block/.style={draw, rounded corners, rectangle, minimum width=2.6cm, minimum height=1.1cm, align=center, font=\small, fill=blue!6},
  monblock/.style={draw, rounded corners, rectangle, minimum width=2.6cm, minimum height=1.1cm, align=center, font=\small, fill=orange!12},
  arrow/.style={-Latex, thick}
]
  \node[block] (gen)   {Packet\\Generator (TB)};
  \node[block, right=1.0cm of gen]  (pcie)  {\texttt{pcie\_model}};
  \node[block, right=1.0cm of pcie] (dma)   {\texttt{dma\_model}};
  \node[block, right=1.0cm of dma]  (mm2s)  {\texttt{axi\_mm2s}\\\texttt{\_adapter}};

  \node[block, below=1.6cm of mm2s]  (fifo)  {\texttt{axis\_fifo}};
  \node[block, left=1.0cm of fifo]   (batch) {\texttt{batching}\\\texttt{\_unit}};
  \node[block, left=1.0cm of batch]  (sink)  {\texttt{stream}\\\texttt{\_sink}};
  \node[monblock, left=1.0cm of sink] (mon)   {\texttt{latency}\\\texttt{\_monitor}};

  \draw[arrow] (gen)  -- node[above,font=\tiny]{gen\_*} (pcie);
  \draw[arrow] (pcie) -- node[above,font=\tiny]{pcie\_m\_*} (dma);
  \draw[arrow] (dma)  -- node[above,font=\tiny]{dma\_m\_*} (mm2s);
  \draw[arrow] (mm2s.south) -- ++(0,-0.55) -| node[pos=0.25,right,font=\tiny]{mm2s\_m\_*} (fifo.north);
  \draw[arrow] (fifo)  -- node[above,font=\tiny]{fifo\_m\_*} (batch);
  \draw[arrow] (batch) -- node[above,font=\tiny]{batch\_m\_*} (sink);
  \draw[arrow] (sink)  -- node[above,font=\tiny]{pkt\_done}   (mon);
  \draw[arrow] (gen.north) -- ++(0,0.55) -| node[pos=0.75,right,font=\tiny]{pkt\_inject} (mon.north);
\end{tikzpicture}
\caption{Pipeline architecture of the PCIe-AXI DMA streaming model. Solid arrows carry AXI4-Stream beats (\texttt{tvalid/tready/tdata/tkeep/tlast/tuser/tid}); the top-to-bottom arrow into \texttt{latency\_monitor} carries only the two scalar injection/completion event signals, not the data stream itself.}
\label{fig:pipeline_overview}
\end{figure}

Every solid arrow in \autoref{fig:pipeline_overview} is a full AXI4-Stream connection (Section~\ref{sec:ch3_axi}): \texttt{tvalid}, \texttt{tready}, \texttt{tdata}, \texttt{tkeep}, \texttt{tlast}, \texttt{tuser}, and \texttt{tid}. Because every stage implements its own independent \texttt{tvalid}/\texttt{tready} handshake, back-pressure asserted anywhere downstream (for example, a full FIFO) propagates correctly upstream through every intervening stage without any stage needing to know about the others -- a direct consequence of the AXI4-Stream composability property noted in \autoref{sec:ch3_axi}.

%----------------------------------------------------------------------------------------

\subsection{Shared Package and Sideband Signalling Convention}
\label{sec:ch4_sideband}

All modules and the testbench import a single SystemVerilog package, \texttt{pcie\_axi\_pkg}, which centralises every width and timing parameter (data width, packet-size field width, PCIe and DMA timing constants, FIFO depth, batch size, and so on) so that a single edit propagates consistently to every stage and to the sweep experiments of \autoref{Chapter7}. \autoref{tab:pkg_params} lists the default parameter values used throughout this thesis unless a specific experiment overrides them.

\begin{table}[htbp]
\centering
\begin{tabular}{@{}llrl@{}}
\toprule
\textbf{Parameter} & \textbf{Meaning} & \textbf{Default} & \textbf{Unit} \\
\midrule
\texttt{DATA\_WIDTH}        & AXI-Stream data bus width         & 64  & bits (8 B/beat) \\
\texttt{PKT\_SIZE\_W}       & Width of the packet-size field     & 13  & bits (0--8191 B) \\
\texttt{PKT\_ID\_W}         & Width of the packet tag field       & 8   & bits (0--255) \\
\texttt{PCIE\_BASE\_LATENCY}   & PCIe fixed per-TLP overhead      & 20  & cycles \\
\texttt{PCIE\_PER\_DW\_LATENCY}& PCIe per-DWORD overhead          & 1   & cycles/DWORD \\
\texttt{DMA\_SETUP\_CYCLES}    & DMA descriptor-fetch overhead    & 8   & cycles \\
\texttt{DMA\_BYTES\_PER\_CYC}  & DMA effective transfer rate      & 8   & bytes/cycle \\
\texttt{FIFO\_DEPTH\_DEFAULT}  & AXI-Stream FIFO depth            & 64  & beats \\
\texttt{BATCH\_COUNT\_DEFAULT} & Packets coalesced per batch      & 4   & packets \\
\texttt{BATCH\_TIMEOUT\_DEFAULT}& Idle cycles before forced flush & 64  & cycles \\
\bottomrule
\end{tabular}
\caption{Default parameters exported by \texttt{pcie\_axi\_pkg} and used throughout this thesis unless stated otherwise.}
\label{tab:pkg_params}
\end{table}

Two AXI4-Stream sideband fields are given a fixed, pipeline-wide meaning so that every stage, and the latency monitor, can recover per-packet metadata without a dedicated side channel:
\begin{itemize}
    \item \texttt{tuser} carries the packet's byte length (constant across every beat of a packet), which every stage that needs to compute a size-dependent delay (\texttt{pcie\_model}, \texttt{dma\_model}) or a size-dependent \texttt{tkeep} mask (\texttt{axi\_mm2s\_adapter}) reads directly rather than recomputing from a beat count.
    \item \texttt{tid} carries the packet's identifying tag, which the latency monitor uses to correlate an injection event with its corresponding completion event, since packets may in general be in flight and out of order with respect to other tagged events.
\end{itemize}

%----------------------------------------------------------------------------------------

\section{\texttt{pcie\_model}: PCIe Transaction-Layer Model}
\label{sec:ch4_pcie_model}

\texttt{pcie\_model} implements Equation~\ref{eq:pcie_latency} as a three-state finite-state machine (receive, delay, send). It buffers an entire incoming packet, computes its PCIe delay from the latched packet size once the packet's final beat (\texttt{tlast}) is seen, holds the packet for exactly that many cycles while de-asserting \texttt{s\_tready} (back-pressuring the packet generator), and then replays the buffered beats to \texttt{dma\_model}. Because the delay is computed once per packet from the packet's total size rather than per beat, this module reproduces the fixed-plus-proportional PCIe cost of \autoref{sec:ch3_pcie} exactly, including the tighter case of the module's internal buffer having to hold up to a full 4~KB packet (512 beats at the default 64-bit bus width). The detailed state-transition behaviour is presented in \autoref{sec:ch5_pcie_fsm}.

\section{\texttt{dma\_model}: DMA Engine Model}
\label{sec:ch4_dma_model}

\texttt{dma\_model} implements Equation~\ref{eq:dma_latency} with an analogous three-state structure (receive, setup, send). After buffering a full packet from \texttt{pcie\_model}, it spends a fixed \texttt{SETUP\_CYCLES} modelling descriptor fetch and doorbell processing (\autoref{sec:ch3_dma}), and then plays the packet out at a rate governed by \texttt{DMA\_BYTES\_PER\_CYC}: at the default rate (equal to the bus width), one beat is produced per cycle with no throttling; at a reduced rate, an inter-beat holdoff models a bandwidth-limited memory read. This throttle exists so the model can, if required, separate DMA-side bandwidth limitations from the fixed setup overhead that is the focus of this thesis. \autoref{sec:ch5_dma_fsm} gives the full state-transition detail.

\section{\texttt{axi\_mm2s\_adapter}: Memory-Mapped-to-Stream Adapter}
\label{sec:ch4_mm2s}

\texttt{axi\_mm2s\_adapter} models the final re-packetisation stage of a real \ac{MM2S} DMA path, in which pre-staged payload bytes are serialised onto an AXI4-Stream bus. Because this thesis's upstream stages already deliver data beat-by-beat, the adapter's only substantive function is combinational: it computes the correct \texttt{tkeep} mask for the final beat of every packet, using the packet's byte length (via \texttt{tuser}) modulo the bus width in bytes, so that a packet whose length is not an exact multiple of the bus width correctly marks only its valid trailing bytes. The adapter adds no pipeline stages and hence no latency of its own -- architecturally correct, since in a real system this logic is folded into the DMA engine itself rather than existing as a separate, registered stage.

\section{\texttt{axis\_fifo}: Parameterised AXI-Stream FIFO}
\label{sec:ch4_fifo}

\texttt{axis\_fifo} is a synchronous, first-word-fall-through (\ac{FWFT}) circular buffer of configurable depth \texttt{DEPTH}, decoupling the (potentially bursty) production rate of the upstream DMA path from the consumption rate of whatever sits downstream. Because its read side is combinational (the head-of-queue entry is presented on \texttt{m\_tdata} as soon as the FIFO is non-empty, with no extra registered pull stage), the FIFO adds zero cycles of latency to a packet that finds it empty on arrival -- latency only appears once the write side stalls because the FIFO is full. This property is what makes the FIFO-depth experiment of \autoref{sec:ch7_fifo_stall} meaningful only under a rate-limited consumer: with an always-ready sink, the FIFO drains every cycle it holds data and can never accumulate occupancy, so depth has no observable effect (\autoref{sec:ch7_fifo_always_ready}). The module also exposes an \texttt{occupancy} output and \texttt{full}/\texttt{empty} flags used directly by the testbench's FIFO-stress instrumentation (\autoref{sec:ch6_testbench}).

\section{\texttt{batching\_unit}: Small-Packet Batching Unit}
\label{sec:ch4_batching_unit}

\texttt{batching\_unit} is the RTL realisation of the overhead-amortisation principle motivated in \autoref{sec:ch2_batching} and formalised in Equation~\ref{eq:batch_inequality}. It is an online, zero-buffering pass-through stage: every beat flows from its slave port to its master port on the same cycle, and the unit modifies only the packet-boundary signalling, not the data itself. While counting inner packets, it suppresses every \texttt{tlast} that would otherwise mark the end of an individual small packet, and instead asserts a single \texttt{tlast} once either \texttt{BATCH\_COUNT} packets have been accumulated or an idle-cycle timeout (\texttt{BATCH\_TIMEOUT}) forces an early flush. The resulting single, larger frame carries the accumulated total byte count on \texttt{tuser} and the tag of its first constituent packet on \texttt{tid}, so that downstream stages -- in particular the latency monitor -- correctly attribute the whole coalesced frame to one measurement rather than to its first inner packet's original size. When disabled (\texttt{enable=0}), the unit degenerates to a zero-latency combinational pass-through, which is what makes the un-batched baseline of \autoref{Chapter7} a fair, single-parameter A/B comparison against the batched configuration. The full batching algorithm and its corner cases are detailed in \autoref{sec:ch5_batching}.

\section{\texttt{stream\_sink}: Terminal Consumer and Completion Signalling}
\label{sec:ch4_stream_sink}

\texttt{stream\_sink} represents the endpoint consumer of the pipeline -- conceptually, a PCIe device receive buffer, a network MAC, or a host-side write target. By default it is permanently ready (\texttt{s\_tready = 1}), so that the datapath's measured latency reflects only the PCIe/DMA/buffering path and is not confounded by consumer-side back-pressure; an optional configurable stall model (accepting only a fraction of cycles, parameterised by \texttt{STALL\_PERIOD}/\texttt{STALL\_READY}) is provided specifically to exercise the FIFO-depth experiment of \autoref{sec:ch7_fifo_stall}, since an always-ready sink never allows the upstream FIFO to accumulate occupancy. On every beat that is both valid and accepted with \texttt{tlast} asserted, the sink emits a one-cycle \texttt{pkt\_done} pulse together with the completing packet's tag (\texttt{done\_id}) and its actual byte count (\texttt{done\_size}, batch-aware -- i.e. the total frame size when the batching unit is active) to the latency monitor.

\section{\texttt{latency\_monitor}: Hardware Latency Measurement}
\label{sec:ch4_latency_monitor}

\texttt{latency\_monitor} is a hardware performance-counter block, not a testbench artefact: it is instantiated inside the design itself and measures latency using only signals available in hardware. It maintains a free-running cycle counter and, indexed by packet tag, a table of injection timestamps and injected sizes. When a \texttt{pkt\_inject} event arrives (asserted by the packet generator at the same cycle the first beat of a packet is presented), the current cycle count and size are recorded against that packet's tag; when the corresponding \texttt{pkt\_done} event arrives from \texttt{stream\_sink}, the elapsed cycle count since injection is computed and registered as that packet's latency, together with running minimum, maximum, average, and throughput statistics reported at the end of simulation. Because injection and completion are both hardware events on the same clock domain, this measurement is immune to the kind of software-side reference-point misalignment that produced a spurious constant offset in an earlier iteration of the testbench (\autoref{sec:ch6_milestones}), and its output is what is reported directly as the \texttt{[RESULT]} rows analysed in \autoref{Chapter7}.

%----------------------------------------------------------------------------------------

\section{Top-Level Integration and Parameterisation}
\label{sec:ch4_top}

\texttt{pcie\_axi\_dma\_top} instantiates and wires together the seven stages of \autoref{fig:pipeline_overview} using a consistent per-stage AXI4-Stream bus naming convention (\texttt{<stage>\_tvalid}, \texttt{<stage>\_tready}, and so on), and re-exports, as module parameters, every configuration knob needed to run the experiments of \autoref{Chapter7} without editing RTL: the PCIe and DMA timing constants, the FIFO depth, the batch count and timeout, a single \texttt{BATCHING\_EN} bit that switches the batching unit between bypass and active modes, and the stream-sink stall parameters. Exposing \texttt{BATCHING\_EN} as a synthesis-time parameter rather than a runtime signal was a deliberate choice: it allows a downstream synthesis tool to prune the entire batching datapath in the baseline configuration, and it guarantees that the baseline and batched configurations of \autoref{sec:ch7_batching} differ in exactly one bit at elaboration time, isolating the effect being measured.
